% Appendix: Formula Derivations
% Complete derivations for Chapter 1-2

\appendix
\chapter{公式推导}

\section{概率与推断（Chapter 1）推导}

\subsection{贝叶斯定理的推导}\label{sec:bayes-theorem-derivation}

\textbf{背景（Background）}：
贝叶斯定理是概率论中最基本的公式之一，它描述了如何根据新的证据更新概率估计。这个推导基于条件概率的定义和基本运算规则。

\textbf{参数定义（Parameter Definitions）}：
\begin{itemize}
  \item $A$：某个事件
  \item $B$：另一个事件（通常代表观测数据或证据）
  \item $p(A)$：事件 $A$ 的先验概率（prior probability）
  \item $p(A|B)$：在事件 $B$ 发生的条件下，事件 $A$ 的后验概率（posterior probability）
  \item $p(B|A)$：在事件 $A$ 发生的条件下，事件 $B$ 的似然（likelihood）
  \item $p(B)$：事件 $B$ 的边际概率（marginal probability）
\end{itemize}

\textbf{已知条件（Given）}：
\begin{itemize}
  \item 条件概率定义：$p(A|B) = \dfrac{p(A \cap B)}{p(B)}$，其中 $p(B) > 0$
  \item 乘法法则：$p(A \cap B) = p(A) \cdot p(B|A) = p(B) \cdot p(A|B)$
\end{itemize}

\textbf{目标（Goal）}：
从条件概率的基本定义推导出贝叶斯定理：
\[
p(A|B) = \frac{p(A) \cdot p(B|A)}{p(B)}
\]

\textbf{推导步骤（Derivation Steps）}：

\textbf{步骤 1}：从条件概率定义出发
\[
p(A|B) = \frac{p(A \cap B)}{p(B)}
\tag*{条件概率定义}

\textbf{步骤 2}：用乘法法则替换联合概率 $p(A \cap B)$
\[
p(A \cap B) = p(A) \cdot p(B|A)
\tag*{乘法法则}

\textbf{步骤 3}：将步骤 2 代入步骤 1
\[
p(A|B) = \frac{p(A) \cdot p(B|A)}{p(B)}
\tag*{代入即得}

\textbf{步骤 4}：对 $p(B)$ 进行边际化（如果 $A_1, A_2, \ldots, A_k$ 是互斥且完备的事件族）
\[
p(B) = \sum_{i=1}^{k} p(A_i) \cdot p(B|A_i)
\tag*{全概率公式}

\textbf{步骤 5}：将步骤 4 代入贝叶斯定理的分母，得到完整形式
\[
p(A_j|B) = \frac{p(A_j) \cdot p(B|A_j)}{\sum_{i=1}^{k} p(A_i) \cdot p(B|A_i)}
\tag*{贝叶斯定理（完整形式）}

\textbf{结论}：贝叶斯定理建立了先验概率与后验概率之间的桥梁——我们通过似然 $p(B|A)$ 将新的观测 $B$ 与先验知识 $p(A)$ 结合起来。

\clearpage

\subsection{先验预测分布与后验预测分布的推导}\label{sec:predictive-distributions-derivation}

\textbf{背景（Background）}：
在贝叶斯统计中，我们不仅要推断当前参数 $\theta$ 的后验分布，还需要预测未来观测 $\tilde{y}$ 的分布。这涉及对未知参数进行积分（边际化）。

\textbf{参数定义（Parameter Definitions）}：
\begin{itemize}
  \item $\theta$：不可观测的参数
  \item $y$：已观测的数据
  \item $\tilde{y}$：未来或未观测的数据
  \item $p(\theta)$：参数的先验分布
  \item $p(y|\theta)$：数据的似然函数
  \item $p(\theta|y)$：参数的后验分布
\end{itemize}

\textbf{已知条件（Given）}：
\begin{itemize}
  \item 贝叶斯定理：$p(\theta|y) \propto p(\theta) \cdot p(y|\theta)$
  \item 边际概率定义：$p(y) = \int p(\theta) \cdot p(y|\theta) \, d\theta$
  \item 条件独立性：给定 $\theta$，未来观测 $\tilde{y}$ 与当前观测 $y$ 条件独立
\end{itemize}

\textbf{目标（Goal）}：
\begin{enumerate}
  \item 推导先验预测分布 $p(y)$
  \item 推导后验预测分布 $p(\tilde{y}|y)$
\end{enumerate}

\textbf{推导步骤（Derivation Steps）}：

\subsubsection*{先验预测分布的推导}

\textbf{步骤 1}：从边际概率的定义出发
\[
p(y) = \int p(\theta) \cdot p(y|\theta) \, d\theta
\tag*{边际概率定义}

\textbf{步骤 2}：解释其含义——这是对所有可能的 $\theta$ 值加权平均，其中权重是先验概率

\textbf{步骤 3}：先验预测分布表示在我们观测数据之前，所有可能观测结果的分布

\clearpage

\subsubsection*{后验预测分布的推导}

\textbf{步骤 1}：从条件概率定义出发
\[
p(\tilde{y}|y) = \int p(\tilde{y}, \theta|y) \, d\theta
\tag*{条件概率的边际化}

\textbf{步骤 2}：将联合分布分解
\[
p(\tilde{y}, \theta|y) = p(\tilde{y}|\theta, y) \cdot p(\theta|y)
\tag*{链式法则}

\textbf{步骤 3}：利用条件独立性假设（给定 $\theta$，$\tilde{y}$ 与 $y$ 独立）
\[
p(\tilde{y}|\theta, y) = p(\tilde{y}|\theta)
\tag*{条件独立性}

\textbf{步骤 4}：代入得
\[
p(\tilde{y}|y) = \int p(\tilde{y}|\theta) \cdot p(\theta|y) \, d\theta
\tag*{后验预测分布}

\textbf{结论}：后验预测分布是\textbf{在参数后验分布上对条件预测进行平均}。这体现了贝叶斯推断的核心思想：不确定性在两个层次上都存在（参数的不确定性和预测的随机性），我们需要对两者都进行边际化。

\clearpage

\subsection{后验优势比的推导}\label{sec:posterior-odds-derivation}

\textbf{背景（Background）}：
优势比（odds ratio）是概率的另一种表示方式，在许多应用场景下比概率更直观。贝叶斯定理在优势比形式下有一个特别简洁的表达。

\textbf{参数定义（Parameter Definitions）}：
\begin{itemize}
  \item $\theta_1$：某个参数值
  \item $\theta_2$：另一个参数值
  \item $p(\theta_1)$：$\theta_1$ 的先验概率
  \item $p(\theta_1|y)$：$\theta_1$ 的后验概率
  \item $p(y|\theta_1)$：似然函数在 $\theta_1$ 处的值
  \item $O(\theta_1,\theta_2) = \dfrac{p(\theta_1)}{p(\theta_2)}$：先验优势比
  \item $O(\theta_1,\theta_2|y) = \dfrac{p(\theta_1|y)}{p(\theta_2|y)}$：后验优势比
\end{itemize}

\textbf{已知条件（Given）}：
贝叶斯定理：
\[
p(\theta|y) = \frac{p(\theta) \cdot p(y|\theta)}{p(y)}
\]

\textbf{目标（Goal）}：
证明后验优势比 = 先验优势比 × 似然比

\textbf{推导步骤（Derivation Steps）}：

\textbf{步骤 1}：写出后验优势比的定义
\[
O(\theta_1,\theta_2|y) = \frac{p(\theta_1|y)}{p(\theta_2|y)}
\tag*{后验优势比定义}

\textbf{步骤 2}：对分子和分母分别应用贝叶斯定理
\[
= \frac{\dfrac{p(\theta_1) \cdot p(y|\theta_1)}{p(y)}}{\dfrac{p(\theta_2) \cdot p(y|\theta_2)}{p(y)}}
\tag*{代入贝叶斯定理}

\textbf{步骤 3}：约去 $p(y)$
\[
= \frac{p(\theta_1) \cdot p(y|\theta_1)}{p(\theta_2) \cdot p(y|\theta_2)}
\tag*{约去边际概率}

\textbf{步骤 4}：重新组合
\[
= \frac{p(\theta_1)}{p(\theta_2)} \times \frac{p(y|\theta_1)}{p(y|\theta_2)}
\tag*{即先验优势比乘以似然比}

\textbf{结论}：这个简洁的关系式告诉我们，\textbf{后验优势比 = 先验优势比 × 似然比}。数据通过似然比来更新优势比——如果 $p(y|\theta_1) > p(y|\theta_2)$，则数据支持 $\theta_1$，后验优势比增大。

\clearpage

\section{单参数模型（Chapter 2）推导}

\subsection{Beta-Binomial 共轭后验分布推导}\label{sec:beta-binomial-posterior-derivation}

\textbf{背景（Background）}：
在二项模型中，如果我们使用 Beta 先验分布来估计成功概率 $\theta$，则后验分布也是 Beta 分布。这种性质称为共轭性，是贝叶斯统计中最重要的概念之一。

\textbf{参数定义（Parameter Definitions）}：
\begin{itemize}
  \item $n$：伯努利试验的总次数
  \item $y$：$n$ 次试验中成功的次数
  \item $\theta$：每次试验的成功概率（未知参数）
  \item $\alpha$：Beta 先验的第一个参数
  \item $\beta$：Beta 先验的第二个参数
  \item $p(y|\theta) = \binom{n}{y} \theta^y (1-\theta)^{n-y}$：二项似然函数
  \item $p(\theta) = \dfrac{\Gamma(\alpha+\beta)}{\Gamma(\alpha)\Gamma(\beta)} \theta^{\alpha-1}(1-\theta)^{\beta-1}$：Beta 先验
\end{itemize}

\textbf{已知条件（Given）}：
\begin{itemize}
  \item Beta 函数定义：$B(\alpha, \beta) = \dfrac{\Gamma(\alpha)\Gamma(\beta)}{\Gamma(\alpha+\beta)}$
  \item 似然函数：$p(y|\theta) \propto \theta^y (1-\theta)^{n-y}$（忽略不依赖 $\theta$ 的常数 $\binom{n}{y}$）
  \item 先验分布：$p(\theta) \propto \theta^{\alpha-1}(1-\theta)^{\beta-1}$
\end{itemize}

\textbf{目标（Goal）}：
证明后验分布 $p(\theta|y)$ 也是 Beta 分布，并确定其后验参数。

\textbf{推导步骤（Derivation Steps）}：

\textbf{步骤 1}：应用贝叶斯定理写出非归一化后验分布
\[
p(\theta|y) \propto p(\theta) \cdot p(y|\theta)
\tag*{贝叶斯定理（非归一化形式）}

\textbf{步骤 2}：代入先验和似然
\[
p(\theta|y) \propto \theta^{\alpha-1}(1-\theta)^{\beta-1} \cdot \theta^y (1-\theta)^{n-y}
\tag*{代入先验和似然}

\textbf{步骤 3}：合并同底数指数
\[
= \theta^{(\alpha-1)+y} \cdot (1-\theta)^{(\beta-1)+(n-y)}
= \theta^{\alpha+y-1} \cdot (1-\theta)^{\beta+n-y-1}
\tag*{合并指数}

\textbf{步骤 4}：识别分布形式
这正是 Beta 分布的核（kernel），即
\[
p(\theta|y) \propto \theta^{\alpha_{\text{post}}-1}(1-\theta)^{\beta_{\text{post}}-1}
\]
其中
\[
\alpha_{\text{post}} = \alpha + y, \quad \beta_{\text{post}} = \beta + n - y
\tag*{后验参数}

\textbf{步骤 5}：写出完整的后验分布
\[
\theta|y \sim \text{Beta}(\alpha + y, \beta + n - y)
\tag*{后验分布}

\textbf{结论}：后验分布与先验分布属于同一形式（Beta 分布），这就是\textbf{共轭性}。先验参数 $\alpha$ 可以解释为 $\alpha - 1$ 次"先验成功"，$\beta$ 解释为 $\beta - 1$ 次"先验失败"。

\clearpage

\subsection{Beta-Binomial 后验均值与方差推导}\label{sec:beta-binomial-moments-derivation}

\textbf{背景（Background）}：
在共轭框架下，我们可以推导出后验分布的均值和方差的闭式表达式。这些表达式展示了先验信息与数据信息如何"折中"。

\textbf{参数定义（Parameter Definitions）}：
\begin{itemize}
  \item $\alpha_{\text{post}} = \alpha + y$：后验第一个参数
  \item $\beta_{\text{post}} = \beta + n - y$：后验第二个参数
  \item $\mathbb{E}(\theta|y)$：后验均值
  \item $\text{Var}(\theta|y)$：后验方差
\end{itemize}

\textbf{已知条件（Given）}：
\begin{itemize}
  \item 后验分布：$\theta|y \sim \text{Beta}(\alpha_{\text{post}}, \beta_{\text{post}})$
  \item Beta 分布均值公式：$\mathbb{E}(X) = \dfrac{\alpha}{\alpha+\beta}$，其中 $X \sim \text{Beta}(\alpha, \beta)$
  \item Beta 分布方差公式：$\text{Var}(X) = \dfrac{\alpha\beta}{(\alpha+\beta)^2(\alpha+\beta+1)}$
\end{itemize}

\textbf{目标（Goal）}：
\begin{enumerate}
  \item 求后验均值 $\mathbb{E}(\theta|y)$
  \item 求后验方差 $\text{Var}(\theta|y)$
\end{enumerate}

\clearpage

\subsubsection*{后验均值推导}

\textbf{步骤 1}：应用 Beta 分布均值公式
\[
\mathbb{E}(\theta|y) = \frac{\alpha_{\text{post}}}{\alpha_{\text{post}} + \beta_{\text{post}}}
\tag*{Beta 分布均值公式}

\textbf{步骤 2}：代入后验参数
\[
= \frac{\alpha + y}{\alpha + y + \beta + n - y}
= \frac{\alpha + y}{\alpha + \beta + n}
\tag*{代入参数}

\textbf{步骤 3}：改写为加权平均形式
\[
= \frac{\alpha + y}{\alpha + \beta + n}
= \frac{\alpha + \beta}{\alpha + \beta + n} \cdot \frac{\alpha}{\alpha + \beta}
  + \frac{n}{\alpha + \beta + n} \cdot \frac{y}{n}
\tag*{加权平均分解}

\textbf{结论 1}：后验均值是先验均值 $\dfrac{\alpha}{\alpha+\beta}$ 和样本比例 $\dfrac{y}{n}$ 的加权平均，权重与样本量和先验强度相关。

\clearpage

\subsubsection*{后验方差推导}

\textbf{步骤 1}：应用 Beta 分布方差公式
\[
\text{Var}(\theta|y) = \frac{\alpha_{\text{post}} \cdot \beta_{\text{post}}}{(\alpha_{\text{post}} + \beta_{\text{post}})^2 (\alpha_{\text{post}} + \beta_{\text{post}} + 1)}
\tag*{Beta 分布方差公式}

\textbf{步骤 2}：代入后验参数
\[
= \frac{(\alpha + y)(\beta + n - y)}{(\alpha + \beta + n)^2 (\alpha + \beta + n + 1)}
\tag*{代入参数}

\textbf{步骤 3}：大样本近似（当 $y$ 和 $n-y$ 较大时）
当 $n$ 很大而 $\alpha, \beta$ 固定时，
\[
\text{Var}(\theta|y) \approx \frac{1}{n} \cdot \frac{y}{n}\left(1 - \frac{y}{n}\right)
\tag*{大样本近似}

\textbf{结论 2}：后验方差以 $1/n$ 的速率趋近于零，表明随着样本量增大，后验分布越来越集中在真值附近。

\clearpage

\subsection{拉普拉斯继承法则推导}\label{sec:laplace-succession-derivation}

\textbf{背景（Background）}：
拉普拉斯在 18 世纪研究"明天太阳会升起"这类问题时，提出了一个著名的法则。在均匀先验下，这个法则预测未来成功的概率为 $\dfrac{y+1}{n+2}$。

\textbf{参数定义（Parameter Definitions）}：
\begin{itemize}
  \item $n$：已完成的试验次数
  \item $y$：成功的次数
  \item $\tilde{y}$：下一次试验的结果（$\tilde{y} = 1$ 表示成功）
  \item $\theta$：成功概率
  \item $p(\theta|y) = \text{Beta}(y+1, n-y+1)$：均匀先验下的后验分布
\end{itemize}

\textbf{已知条件（Given）}：
\begin{itemize}
  \item 均匀先验：$\theta \sim \text{Uniform}(0,1) = \text{Beta}(1,1)$
  \item 后验分布：$\theta|y \sim \text{Beta}(y+1, n-y+1)$
  \item 后验预测分布：$p(\tilde{y}=1|y) = \int_0^1 \theta \cdot p(\theta|y) \, d\theta$
\end{itemize}

\textbf{目标（Goal）}：
求下一次试验成功的概率 $p(\tilde{y}=1|y)$。

\textbf{推导步骤（Derivation Steps）}：

\textbf{步骤 1}：写出后验预测分布的定义
\[
p(\tilde{y}=1|y) = \int_0^1 p(\tilde{y}=1|\theta) \cdot p(\theta|y) \, d\theta
\tag*{后验预测分布定义}

\textbf{步骤 2}：代入 $p(\tilde{y}=1|\theta) = \theta$（给定 $\theta$，成功的条件概率就是 $\theta$）
\[
= \int_0^1 \theta \cdot p(\theta|y) \, d\theta
= \mathbb{E}(\theta|y)
\tag*{代入条件概率}

\textbf{步骤 3}：识别这是后验均值
\[
= \mathbb{E}(\theta|y)
\tag*{后验均值}

\textbf{步骤 4}：应用 Beta 分布均值公式
\[
= \frac{y+1}{(y+1) + (n-y+1)}
= \frac{y+1}{n+2}
\tag*{Beta 均值公式}

\textbf{结论}：拉普拉斯继承法则表明，在均匀先验下，已知 $n$ 次试验中 $y$ 次成功，下一次成功的概率是 $\dfrac{y+1}{n+2}$。这推广了"瑞士表匠"问题的直觉——即使观测到 $n$ 次成功，概率也不是 1（因为 $\theta$ 本身可能是变化的）。

\clearpage

\subsection{总期望与总方差公式推导}\label{sec:law-of-total-expectation-variance-derivation}

\textbf{背景（Background）}：
全期望公式和全方差公式是概率论中的两个基本恒等式，它们描述了如何通过条件分布来计算无条件矩。

\textbf{参数定义（Parameter Definitions）}：
\begin{itemize}
  \item $\theta$：随机变量（参数）
  \item $y$：随机变量（数据）
  \item $\mathbb{E}(\theta)$：$\theta$ 的无条件（先验）均值
  \item $\text{Var}(\theta)$：$\theta$ 的无条件（先验）方差
  \item $\mathbb{E}(\theta|y)$：给定 $y$ 后 $\theta$ 的条件均值
  \item $\text{Var}(\theta|y)$：给定 $y$ 后 $\theta$ 的条件方差
\end{itemize}

\textbf{已知条件（Given）}：
\begin{itemize}
  \item 条件期望定义：$\mathbb{E}(\theta|y)$ 是 $y$ 的函数
  \item 迭代期望法则：$\mathbb{E}[\mathbb{E}(\theta|y)] = \mathbb{E}(\theta)$
\end{itemize}

\textbf{目标（Goal）}：
\begin{enumerate}
  \item 证明全期望公式：$\mathbb{E}(\theta) = \mathbb{E}[\mathbb{E}(\theta|y)]$
  \item 证明全方差公式：$\text{Var}(\theta) = \mathbb{E}[\text{Var}(\theta|y)] + \text{Var}[\mathbb{E}(\theta|y)]$
\end{enumerate}

\clearpage

\subsubsection*{全期望公式推导}

\textbf{步骤 1}：从条件期望的定义出发
\[
\mathbb{E}[\mathbb{E}(\theta|y)] = \int \mathbb{E}(\theta|y) \cdot p(y) \, dy
\tag*{外层期望定义}

\textbf{步骤 2}：代入条件期望 $\mathbb{E}(\theta|y) = \int \theta \cdot p(\theta|y) \, d\theta$
\[
= \int \left( \int \theta \cdot p(\theta|y) \, d\theta \right) p(y) \, dy
\tag*{代入条件期望定义}

\textbf{步骤 3}：交换积分次序（由富比尼定理保证）
\[
= \int \theta \left( \int p(\theta|y) p(y) \, dy \right) d\theta
\tag*{交换积分次序}

\textbf{步骤 4}：识别内层积分 $p(\theta) = \int p(\theta|y) p(y) dy$（全概率）
\[
= \int \theta \cdot p(\theta) \, d\theta
= \mathbb{E}(\theta)
\tag*{边际分布定义}

\textbf{结论 1}：全期望公式表明，$\theta$ 的先验均值等于在所有可能数据上后验均值的加权平均。

\clearpage

\subsubsection*{全方差公式推导}

\textbf{步骤 1}：从方差定义出发
\[
\text{Var}(\theta) = \mathbb{E}(\theta^2) - [\mathbb{E}(\theta)]^2
\tag*{方差定义}

\textbf{步骤 2}：对 $\mathbb{E}(\theta^2)$ 应用全期望公式
\[
\mathbb{E}(\theta^2) = \mathbb{E}[\mathbb{E}(\theta^2|y)]
\tag*{全期望公式应用于 $\theta^2$}

\textbf{步骤 3}：在给定 $y$ 的条件下，方差可以写为
\[
\text{Var}(\theta|y) = \mathbb{E}(\theta^2|y) - [\mathbb{E}(\theta|y)]^2
\tag*{条件方差定义}

\textbf{步骤 4}：解出 $\mathbb{E}(\theta^2|y)$
\[
\mathbb{E}(\theta^2|y) = \text{Var}(\theta|y) + [\mathbb{E}(\theta|y)]^2
\tag*{移项}

\textbf{步骤 5}：代入步骤 2
\[
\mathbb{E}(\theta^2) = \mathbb{E}[\text{Var}(\theta|y)] + \mathbb{E}[[\mathbb{E}(\theta|y)]^2]
\tag*{代入}

\textbf{步骤 6}：对 $[\mathbb{E}(\theta|y)]^2$ 应用方差分解
\[
\mathbb{E}[[\mathbb{E}(\theta|y)]^2] = \text{Var}[\mathbb{E}(\theta|y)] + \{\mathbb{E}[\mathbb{E}(\theta|y)]\}^2
\tag*{方差与均值平方的关系}

\textbf{步骤 7}：代入全期望公式的结果 $\mathbb{E}[\mathbb{E}(\theta|y)] = \mathbb{E}(\theta)$
\[
= \text{Var}[\mathbb{E}(\theta|y)] + [\mathbb{E}(\theta)]^2
\tag*{简化}

\textbf{步骤 8}：代回方差的定义
\[
\text{Var}(\theta) = \mathbb{E}(\theta^2) - [\mathbb{E}(\theta)]^2
= \mathbb{E}[\text{Var}(\theta|y)] + \text{Var}[\mathbb{E}(\theta|y)]
\tag*{最终结果}

\textbf{结论 2}：全方差公式表明，\textbf{总方差 = 条件方差的均值 + 条件均值的方差}。这在贝叶斯语境下的解释是：参数的总不确定性来自两个方面——（1）给定数据后参数仍有的剩余不确定性（条件方差的均值），以及（2）后验均值随数据变化的不确定性（条件均值的方差）。

\clearpage

\subsection{正态均值估计的后验分布推导}\label{sec:normal-mean-posterior-derivation}

\textbf{背景（Background）}：
当我们用正态似然估计正态均值时，如果使用正态先验，则后验分布也是正态分布。这个推导展示了共轭先验在正态模型中的应用。

\textbf{参数定义（Parameter Definitions）}：
\begin{itemize}
  \item $y \sim \text{Normal}(\theta, \sigma^2)$：观测数据（已知方差 $\sigma^2$）
  \item $\theta \sim \text{Normal}(\mu_0, \tau_0^2)$：正态先验
  \item $\mu_0$：先验均值
  \item $\tau_0^2$：先验方差
  \item $\sigma^2$：已知的观测方差
  \item $\mu_1, \tau_1^2$：后验均值和方差
\end{itemize}

\textbf{已知条件（Given）}：
\begin{itemize}
  \item 似然函数：$p(y|\theta) \propto \exp\left(-\frac{(y-\theta)^2}{2\sigma^2}\right)$
  \item 先验密度：$p(\theta) \propto \exp\left(-\frac{(\theta-\mu_0)^2}{2\tau_0^2}\right)$
  \item 配方法（completing the square）技巧
\end{itemize}

\textbf{目标（Goal）}：
证明后验分布 $\theta|y \sim \text{Normal}(\mu_1, \tau_1^2)$，并求出 $\mu_1$ 和 $\tau_1^2$。

\textbf{推导步骤（Derivation Steps）}：

\textbf{步骤 1}：写出非归一化后验分布的对数
\[
\log p(\theta|y) \propto \log p(\theta) + \log p(y|\theta)
\tag*{对数后验（忽略常数）}

\textbf{步骤 2}：代入先验和似然的对数
\[
= -\frac{(\theta-\mu_0)^2}{2\tau_0^2} - \frac{(y-\theta)^2}{2\sigma^2}
\tag*{代入}

\textbf{步骤 3}：展开平方项
\[
= -\frac{1}{2}\left[\frac{(\theta^2 - 2\mu_0\theta + \mu_0^2)}{\tau_0^2} + \frac{(\theta^2 - 2y\theta + y^2)}{\sigma^2}\right]
\tag*{展开}

\textbf{步骤 4}：合并 $\theta^2$ 项的系数
\[
= -\frac{1}{2}\left[\theta^2\left(\frac{1}{\tau_0^2} + \frac{1}{\sigma^2}\right) - 2\theta\left(\frac{\mu_0}{\tau_0^2} + \frac{y}{\sigma^2}\right) + \text{常数}\right]
\tag*{合并}

\textbf{步骤 5}：识别后验精度（方差的倒数）
\[
\frac{1}{\tau_1^2} = \frac{1}{\tau_0^2} + \frac{1}{\sigma^2}
\tag*{后验精度}

\textbf{步骤 6}：由配方法，后验均值满足
\[
\theta \cdot \frac{1}{\tau_1^2} = \frac{\mu_0}{\tau_0^2} + \frac{y}{\sigma^2}
\tag*{配方法条件}

\textbf{步骤 7}：解出后验均值
\[
\mu_1 = \frac{\mu_0/\tau_0^2 + y/\sigma^2}{1/\tau_0^2 + 1/\sigma^2}
\tag*{后验均值}

\textbf{步骤 8}：写出完整的后验分布
\[
\theta|y \sim \text{Normal}\left(\mu_1, \tau_1^2\right)
\tag*{后验分布}

\textbf{结论}：
\begin{enumerate}
  \item \textbf{精度加法}：后验精度 = 先验精度 + 数据精度
  \item \textbf{后验均值}：是先验均值 $\mu_0$ 和观测值 $y$ 的加权平均，权重与精度成正比
\end{enumerate}

\clearpage

\subsection{Jeffreys 先验与 Fisher 信息推导}\label{sec:jeffreys-prior-derivation}

\textbf{背景（Background）}：
Jeffreys 提出了一种构建非信息先验的系统方法，基于 Fisher 信息的不变性原则。这个先验在参数的单调变换下保持不变。

\textbf{参数定义（Parameter Definitions）}：
\begin{itemize}
  \item $\theta$：参数
  \item $p(y|\theta)$：似然函数
  \item $J(\theta)$：Fisher 信息
  \item $L(\theta) = \log p(y|\theta)$：对数似然
  \item $I(\theta) = -\mathbb{E}\left[\frac{\partial^2 \log p(y|\theta)}{\partial \theta^2}\right]$：Fisher 信息（另一种定义）
  \item $\phi = h(\theta)$：参数的单调变换
\end{itemize}

\textbf{已知条件（Given）}：
\begin{itemize}
  \item Jeffreys 原则：先验应在参数的单调变换下保持不变
  \item Fisher 信息定义：$J(\theta) = -\mathbb{E}\left[\frac{\partial^2 \log p(y|\theta)}{\partial \theta^2}\right]$
  \item 变换公式：$p(\phi) = p(\theta) \cdot \left|\frac{d\theta}{d\phi}\right|$
\end{itemize}

\textbf{目标（Goal）}：
推导 Jeffreys 先验公式并验证其在变换下的不变性。

\textbf{推导步骤（Derivation Steps）}：

\subsubsection*{Jeffreys 先验公式推导}

\textbf{步骤 1}：Jeffreys 提出先验正比于 Fisher 信息的平方根
\[
p(\theta) \propto [J(\theta)]^{1/2}
\tag*{Jeffreys 先验}

\textbf{步骤 2}：验证在变换 $\phi = h(\theta)$ 下的不变性

在新参数 $\phi$ 下，Fisher 信息为
\[
J(\phi) = -\mathbb{E}\left[\frac{\partial^2 \log p(y|\phi)}{\partial \phi^2}\right]
\tag*{Fisher 信息在新参数下}

\textbf{步骤 3}：利用链式法则
\[
\frac{\partial \log p(y|\theta)}{\partial \phi}
= \frac{\partial \log p(y|\theta)}{\partial \theta} \cdot \frac{d\theta}{d\phi}
\tag*{链式法则}

\textbf{步骤 4}：计算二阶导数
\[
\frac{\partial^2 \log p(y|\phi)}{\partial \phi^2}
= \frac{\partial}{\partial \phi}\left(\frac{\partial \log p(y|\theta)}{\partial \theta} \cdot \frac{d\theta}{d\phi}\right)
= \frac{\partial^2 \log p(y|\theta)}{\partial \theta^2} \left(\frac{d\theta}{d\phi}\right)^2 + \frac{\partial \log p(y|\theta)}{\partial \theta} \cdot \frac{d^2\theta}{d\phi^2}
\tag*{求导}

\textbf{步骤 5}：取期望，利用 $\mathbb{E}\left[\frac{\partial \log p(y|\theta)}{\partial \theta}\right] = 0$（得分函数的期望为零）
\[
J(\phi) = \left(\frac{d\theta}{d\phi}\right)^2 J(\theta)
\tag*{Fisher 信息的变换公式}

\textbf{步骤 6}：因此
\[
[J(\phi)]^{1/2} = \left|\frac{d\theta}{d\phi}\right| [J(\theta)]^{1/2}
\tag*{平方根}

\textbf{步骤 7}：新先验为
\[
p(\phi) \propto [J(\phi)]^{1/2} \propto [J(\theta)]^{1/2} \left|\frac{d\theta}{d\phi}\right|
\tag*{变换后的 Jeffreys 先验}

\textbf{结论}：这正是概率密度在变量变换下的标准变换公式，因此 Jeffreys 先验在单调变换下是不变的。

\clearpage

\subsubsection*{二项模型的 Jeffreys 先验推导}

\textbf{步骤 1}：写出二项模型的对数似然
\[
\log p(y|\theta) = \log \binom{n}{y} + y\log\theta + (n-y)\log(1-\theta)
\tag*{对数似然}

\textbf{步骤 2}：计算一阶导数（得分函数）
\[
\frac{\partial}{\partial\theta} \log p(y|\theta) = \frac{y}{\theta} - \frac{n-y}{1-\theta}
\tag*{一阶导数}

\textbf{步骤 3}：计算二阶导数
\[
\frac{\partial^2}{\partial\theta^2} \log p(y|\theta) = -\frac{y}{\theta^2} - \frac{n-y}{(1-\theta)^2}
\tag*{二阶导数}

\textbf{步骤 4}：取负期望（注意 $\mathbb{E}(Y) = n\theta$）
\[
J(\theta) = -\mathbb{E}\left[\frac{\partial^2}{\partial\theta^2} \log p(y|\theta)\right]
= \frac{n\theta}{\theta^2} + \frac{n(1-\theta)}{(1-\theta)^2}
= \frac{n}{\theta(1-\theta)}
\tag*{Fisher 信息}

\textbf{步骤 5}：应用 Jeffreys 公式
\[
p(\theta) \propto \left[\frac{n}{\theta(1-\theta)}\right]^{1/2}
\propto \theta^{-1/2}(1-\theta)^{-1/2}
\tag*{Jeffreys 先验}

\textbf{步骤 6}：识别为 Beta 分布
\[
p(\theta) \propto \theta^{(1/2)-1}(1-\theta)^{(1/2)-1}
\Rightarrow \theta \sim \text{Beta}\left(\frac{1}{2}, \frac{1}{2}\right)
\tag*{Beta(1/2, 1/2)}

\textbf{结论}：二项模型的 Jeffreys 先验是 $\text{Beta}(1/2, 1/2)$，这是一个弱信息先验，比均匀先验 $\text{Beta}(1,1)$ 在边界处更轻。

\clearpage

\subsection{后验预测分布的均值与方差推导}\label{sec:predictive-mean-variance-derivation}

\textbf{背景（Background）}：
后验预测分布 $\tilde{y}|y$ 的均值和方差有两个组成部分——一个来自模型本身的随机性，另一个来自参数的后验不确定性。

\textbf{参数定义（Parameter Definitions）}：
\begin{itemize}
  \item $\tilde{y}$：未来观测
  \item $y$：当前观测
  \item $\theta$：参数
  \item $\mathbb{E}(\tilde{y}|y)$：后验预测均值
  \item $\text{Var}(\tilde{y}|y)$：后验预测方差
\end{itemize}

\textbf{已知条件（Given）}：
\begin{itemize}
  \item $\tilde{y}|\theta \sim \text{Normal}(\theta, \sigma^2)$
  \item $\theta|y \sim \text{Normal}(\mu_1, \tau_1^2)$
  \item 全期望公式：$\mathbb{E}(\tilde{y}|y) = \mathbb{E}[\mathbb{E}(\tilde{y}|\theta, y)]$
  \item 全方差公式：$\text{Var}(\tilde{y}|y) = \mathbb{E}[\text{Var}(\tilde{y}|\theta, y)] + \text{Var}[\mathbb{E}(\tilde{y}|\theta, y)]$
\end{itemize}

\textbf{目标（Goal）}：
求后验预测分布 $\tilde{y}|y$ 的均值和方差。

\clearpage

\subsubsection*{后验预测均值推导}

\textbf{步骤 1}：应用全期望公式
\[
\mathbb{E}(\tilde{y}|y) = \mathbb{E}[\mathbb{E}(\tilde{y}|\theta, y)]
\tag*{全期望公式}

\textbf{步骤 2}：给定 $\theta$，$\tilde{y}$ 与 $y$ 条件独立，且 $\mathbb{E}(\tilde{y}|\theta) = \theta$
\[
= \mathbb{E}(\theta|y)
= \mu_1
\tag*{后验均值}

\textbf{结论 1}：后验预测均值等于参数的后验均值。

\clearpage

\subsubsection*{后验预测方差推导}

\textbf{步骤 1}：应用全方差公式
\[
\text{Var}(\tilde{y}|y) = \mathbb{E}[\text{Var}(\tilde{y}|\theta, y)] + \text{Var}[\mathbb{E}(\tilde{y}|\theta, y)]
\tag*{全方差公式}

\textbf{步骤 2}：第一项——给定 $\theta$ 下的方差是 $\sigma^2$（模型方差）
\[
\mathbb{E}[\text{Var}(\tilde{y}|\theta, y)] = \mathbb{E}(\sigma^2|y) = \sigma^2
\tag*{模型方差}

\textbf{步骤 3}：第二项——条件均值的方差
\[
\text{Var}[\mathbb{E}(\tilde{y}|\theta, y)] = \text{Var}(\theta|y) = \tau_1^2
\tag*{参数不确定性}

\textbf{步骤 4}：合并
\[
\text{Var}(\tilde{y}|y) = \sigma^2 + \tau_1^2
\tag*{后验预测方差}

\textbf{结论 2}：后验预测方差 = \textbf{模型方差} + \textbf{参数后验方差}。这体现了贝叶斯推断的核心：预测的不确定性来自两个方面——（1）模型本身的随机性（$\sigma^2$），以及（2）参数估计的不确定性（$\tau_1^2$）。

% ================================================
% Chapter 11 推导：马尔可夫链模拟基础
% =============================================

\section{拒绝抽样的正确性证明}\label{sec:derivation-rejection-sampling}

\textbf{背景（Background）}：
拒绝抽样通过随机接受机制从目标分布中抽取样本。我们需要证明被接受的样本确实服从目标分布 $p(\theta)$。

\textbf{参数定义（Parameter Definitions）}：
\begin{itemize}
  \item $p(\theta)$：目标（未归一化）分布
  \item $g(\theta)$：提议分布
  \item $M$：覆盖常数，满足 $p(\theta) \leq M \cdot g(\theta)$
  \item $\theta^*$：候选样本
  \item $u^*$：从 $\mathrm{Uniform}(0,1)$ 抽取的随机数
\end{itemize}

\textbf{目标（Goal）}：
证明被接受的样本 $\theta_{\text{accept}}$ 服从目标分布 $p(\theta)$。

\textbf{推导步骤（Derivation Steps）}：

\textbf{步骤 1}：写出接受的联合概率
\[
\mathbb{P}(\text{接受} \cap \theta^* \in d\theta)
= \mathbb{P}(\theta^* \in d\theta) \cdot \mathbb{P}(\text{接受} | \theta^*)
= g(\theta^*) d\theta \cdot \frac{p(\theta^*)}{M g(\theta^*)}
= \frac{p(\theta^*)}{M} d\theta
\tag*{接受联合概率}
\]

\textbf{步骤 2}：计算总接受概率（归一化常数）
\[
\mathbb{P}(\text{接受}) = \int \frac{p(\theta^*)}{M} d\theta = \frac{1}{M}
\tag*{总接受概率}
\]

\textbf{步骤 3}：计算接受样本的条件分布
\[
\mathbb{P}(\theta_{\text{accept}} \in d\theta | \text{接受})
= \frac{\mathbb{P}(\text{接受} \cap \theta^* \in d\theta)}{\mathbb{P}(\text{接受})}
= \frac{p(\theta^*)/M}{1/M} d\theta = p(\theta^*) d\theta
\tag*{条件分布 = 目标分布}
\]

\textbf{结论}：被接受的样本的条件分布正是目标分布 $p(\theta)$。证明完毕。

\clearpage

\subsection{拒绝抽样期望接受率推导}\label{sec:derivation-rejection-acceptance}

\textbf{背景（Background）}：
拒绝抽样的期望接受率是衡量算法效率的关键指标。

\textbf{目标（Goal）}：
推导期望接受率的公式。

\textbf{推导步骤（Derivation Steps）}：

\textbf{步骤 1}：期望接受率定义为
\[
\text{accept rate} = \int \mathbb{P}(\text{接受} | \theta) \cdot g(\theta) \, d\theta
\tag*{期望定义}
\]

\textbf{步骤 2}：代入接受概率
\[
\mathbb{P}(\text{接受} | \theta) = \min\!\left(1, \frac{p(\theta)}{M g(\theta)}\right)
\tag*{接受概率}
\]

\textbf{步骤 3}：于是
\[
\text{accept rate} = \int \min\!\left(\frac{p(\theta)}{M g(\theta)}, 1\right) \cdot g(\theta) \, d\theta
= \frac{1}{M} \int \min\!\bigl(p(\theta), M g(\theta)\bigr) \, d\theta
\tag*{积分结果}
\]

\textbf{步骤 4}：当覆盖条件 $M g(\theta) \geq p(\theta)$ 对所有 $\theta$ 成立时
\[
\text{accept rate} = \frac{1}{M} \int p(\theta) \, d\theta = \frac{1}{M}
\tag*{简化为 $1/M$}
\]

\textbf{结论}：期望接受率与 $M$ 成反比，选择尽可能小的 $M$ 可提高接受率。

\clearpage

\subsection{细致平衡条件推导}\label{sec:derivation-detailed-balance}

\textbf{背景（Background）}：
细致平衡是构造满足特定平稳分布的马尔可夫链的关键条件。

\textbf{目标（Goal）}：
证明如果转移核 $P$ 满足细致平衡，则其平稳分布就是 $\pi$。

\textbf{推导步骤（Derivation Steps）}：

\textbf{步骤 1}：写出细致平衡条件
\[
P(\theta' | \theta) \cdot \pi(\theta) = P(\theta | \theta') \cdot \pi(\theta')
\tag*{细致平衡}
\]

\textbf{步骤 2}：验证 $\pi$ 是平稳分布
\[
\int P(\theta' | \theta) \cdot \pi(\theta) \, d\theta
= \int P(\theta | \theta') \cdot \pi(\theta') \, d\theta
= \pi(\theta') \int P(\theta | \theta') \, d\theta
= \pi(\theta') \cdot 1
= \pi(\theta')
\tag*{平稳分布验证}
\]

其中第三步用到 $\int P(\theta | \theta') d\theta = 1$（转移核的归一化性）。

\textbf{结论}：细致平衡是平稳分布的充分条件，但不是必要条件。

\clearpage

\subsection{Metropolis 算法细致平衡证明}\label{sec:derivation-metropolis}

\textbf{背景（Background）}：
Metropolis 算法的提议分布是对称的，我们需要证明其满足细致平衡。

\textbf{目标（Goal）}：
证明 Metropolis 算法的转移核以目标分布 $p(\theta)$ 为平稳分布。

\textbf{推导步骤（Derivation Steps）}：

\textbf{步骤 1}：Metropolis 转移核
\[
P(\theta' | \theta) = q(\theta' | \theta) \cdot \min\!\left(1, \frac{p(\theta')}{p(\theta)}\right)
\]
其中 $q(\theta' | \theta) = q(\theta | \theta')$（对称提议分布）。

\textbf{步骤 2}：验证细致平衡
\[
P(\theta' | \theta) \cdot p(\theta)
= q(\theta' | \theta) \cdot \min(1, \frac{p(\theta')}{p(\theta)}) \cdot p(\theta)
\]

分两种情况讨论：

\textbf{情况 1}：$p(\theta') \geq p(\theta)$
\[
\text{右边} = q(\theta' | \theta) \cdot 1 \cdot p(\theta)
= q(\theta' | \theta) \cdot p(\theta')
\]

\textbf{情况 2}：$p(\theta') < p(\theta)$
\[
\text{右边} = q(\theta' | \theta) \cdot \frac{p(\theta')}{p(\theta)} \cdot p(\theta)
= q(\theta' | \theta) \cdot p(\theta')
\]

由于 $q(\theta' | \theta) = q(\theta | \theta')$，两种情况都满足
\[
P(\theta' | \theta) \cdot p(\theta) = P(\theta | \theta') \cdot p(\theta')
\tag*{细致平衡成立}
\]

\textbf{结论}：Metropolis 算法的平稳分布正是目标分布 $p(\theta)$。证明完毕。

\clearpage

\subsection{Gibbs 抽样接受率推导}\label{sec:derivation-gibbs-acceptance}

\textbf{背景（Background）}：
Gibbs 抽样是 Metropolis-Hastings 的特例，其接受率恒为 1。我们来推导这个结论。

\textbf{参数定义（Parameter Definitions）}：
\begin{itemize}
  \item $\theta = (\theta_1, \ldots, \theta_d)$：完整参数向量
  \item $\theta_{-j} = (\theta_1, \ldots, \theta_{j-1}, \theta_{j+1}, \ldots, \theta_d)$：除 $\theta_j$ 外的所有参数
  \item $p^*(\theta)$：未归一化的目标分布
  \item $q(\theta_j' | \theta_{-j}) = p(\theta_j' | \theta_{-j}, y)$：提议分布（条件分布）
\end{itemize}

\textbf{目标（Goal）}：
证明 Gibbs 抽样的接受率为 1。

\textbf{推导步骤（Derivation Steps）}：

\textbf{步骤 1}：写出 MH 接受比
\[
r = \frac{p^*(\theta_j', \theta_{-j})}{p^*(\theta_j, \theta_{-j})}
\cdot \frac{q(\theta_j | \theta_{-j})}{q(\theta_j' | \theta_{-j})}
\tag*{MH 接受比}
\]

\textbf{步骤 2}：代入 Gibbs 的提议分布
\[
q(\theta_j | \theta_{-j}) = p(\theta_j | \theta_{-j}, y)
\]
\[
q(\theta_j' | \theta_{-j}) = p(\theta_j' | \theta_{-j}, y)
\tag*{提议分布 = 条件分布}
\]

\textbf{步骤 3}：利用链式法则分解联合分布
\[
p^*(\theta_j', \theta_{-j}) = p^*(\theta_{-j}) \cdot p^*(\theta_j' | \theta_{-j})
\]
\[
p^*(\theta_j, \theta_{-j}) = p^*(\theta_{-j}) \cdot p^*(\theta_j | \theta_{-j})
\tag*{链式法则}
\]

\textbf{步骤 4}：代入接受比
\[
r = \frac{p^*(\theta_{-j}) \cdot p^*(\theta_j' | \theta_{-j})}{p^*(\theta_{-j}) \cdot p^*(\theta_j | \theta_{-j})}
\cdot \frac{p(\theta_j | \theta_{-j}, y)}{p(\theta_j' | \theta_{-j}, y)}
= \frac{p^*(\theta_j' | \theta_{-j})}{p^*(\theta_j | \theta_{-j})}
\cdot \frac{p(\theta_j | \theta_{-j}, y)}{p(\theta_j' | \theta_{-j}, y)}
\tag*{代入}
\]

\textbf{步骤 5}：注意 $p^*(\theta_j | \theta_{-j}, y) \propto p(\theta_j | \theta_{-j}, y)$，因为 $p^*(\theta_{-j})$ 不依赖于 $\theta_j$。于是两个因子互为倒数，
\[
r = 1
\tag*{接受率为 1}
\]

\textbf{结论}：Gibbs 抽样的提议分布恰好是条件分布，因此接受率恒为 1，无需拒绝。

\clearpage

\subsection{分层模型条件分布推导}\label{sec:derivation-hierarchical-conditionals}

\textbf{背景（Background）}：
在分层模型 \cref{eq:hierarchical-model} 中，我们需要推导各个参数的条件后验分布。

\textbf{模型回顾}：
\[
y_j | \theta_j \sim \mathrm{N}(\theta_j, \sigma^2), \quad \theta_j | \mu, \tau \sim \mathrm{N}(\mu, \tau^2)
\]
其中 $\sigma^2$ 已知。

\textbf{目标（Goal）}：
推导 $p(\theta_j | \mu, \tau, y)$、$p(\mu | \theta, \tau, y)$ 和 $p(\tau | \theta, \mu, y)$。

\clearpage

\subsubsection*{$p(\theta_j | \mu, \tau, y)$ 推导}

\textbf{步骤 1}：写出条件分布的核
\[
p(\theta_j | \mu, \tau, y) \propto p(y_j | \theta_j) \cdot p(\theta_j | \mu, \tau)
\propto \exp\!\left(-\frac{(y_j - \theta_j)^2}{2\sigma^2}\right)
\cdot \exp\!\left(-\frac{(\theta_j - \mu)^2}{2\tau^2}\right)
\tag*{核}
\]

\textbf{步骤 2}：合并指数
\[
= \exp\!\left(-\frac{(y_j - \theta_j)^2}{2\sigma^2} - \frac{(\theta_j - \mu)^2}{2\tau^2}\right)
\tag*{合并}
\]

\textbf{步骤 3}：配方得到正态分布形式
\[
= \exp\!\left(-\frac{1}{2}\left[\frac{1}{\sigma^2} + \frac{1}{\tau^2}\right]
\left(\theta_j - \frac{y_j/\sigma^2 + \mu/\tau^2}{1/\sigma^2 + 1/\tau^2}\right)^2\right)
\tag*{配方}
\]

因此
\[
\theta_j | \mu, \tau, y \sim \mathrm{N}\!\left(\frac{y_j/\sigma^2 + \mu/\tau^2}{1/\sigma^2 + 1/\tau^2},
\left(\frac{1}{\sigma^2} + \frac{1}{\tau^2}\right)^{-1}\right)
\tag*{条件后验}
\]

\clearpage

\subsubsection*{$p(\mu | \theta, \tau, y)$ 推导}

\textbf{步骤 1}：$\mu$ 的条件分布（似然部分只依赖于 $\theta_j$，与 $\mu$ 通过 $\theta_j | \mu, \tau$ 连接）
\[
p(\mu | \theta, \tau) \propto p(\theta | \mu, \tau) = \prod_{j=1}^J \exp\!\left(-\frac{(\theta_j - \mu)^2}{2\tau^2}\right)
\propto \exp\!\left(-\frac{\sum_j (\theta_j - \mu)^2}{2\tau^2}\right)
\tag*{先验部分}
\]

\textbf{步骤 2}：展开平方
\[
\sum_j (\theta_j - \mu)^2 = \sum_j (\theta_j - \bar{\theta} + \bar{\theta} - \mu)^2
= \sum_j (\theta_j - \bar{\theta})^2 + J(\bar{\theta} - \mu)^2
\tag*{平方展开}
\]

\textbf{步骤 3}：识别为正态核
\[
p(\mu | \theta, \tau) \propto \exp\!\left(-\frac{J(\bar{\theta} - \mu)^2}{2\tau^2}\right)
\tag*{正态核}
\]

因此
\[
\mu | \theta, \tau \sim \mathrm{N}\!\left(\bar{\theta}, \frac{\tau^2}{J}\right)
\tag*{条件后验}
\]

其中 $\bar{\theta} = \frac{1}{J}\sum_j \theta_j$。

\clearpage

\subsubsection*{$p(\tau | \theta, \mu, y)$ 的非标准形式}

\textbf{说明}：
\[
p(\tau | \theta, \mu, y) \propto p(\theta | \mu, \tau) \cdot p(\tau)
= \prod_{j=1}^J \frac{1}{\tau} \exp\!\left(-\frac{(\theta_j - \mu)^2}{2\tau^2}\right) \cdot p(\tau)
\tag*{非标准形式}
\]

这个分布不是标准形式（不是常见的正态、Gamma 等分布），因此需要使用 Metropolis-Hastings 步骤来抽样。

\textbf{结论}：在 Gibbs 抽样框架下，对 $\tau$ 需要使用 MH 步骤，这正是"混合 MCMC"的典型例子。

% ================================================
% 问答与注解
% ================================================
% Appendix: Questions and Answers
% User annotations and interactive Q&A records

\section{阅读过程中的问题与解答}

% === 用户问答记录 ===
% 本节记录阅读过程中的问答内容

\section{什么是分布的核}\label{sec:qa:DistributionKernel}

\subsection{问题}
分布的核是什么概念？

\subsection{回答}
分布的\textbf{核（kernel）}是指概率密度/质量函数中\textbf{依赖于参数的核心部分}，不包括归一化常数。

\textbf{举例说明}：

Beta-Binomial 共轭中，后验分布：
\[
p(\theta|y) \propto \theta^{\alpha+y-1}(1-\theta)^{\beta+n-y-1}
\]

其中 $\theta^{\alpha+y-1}(1-\theta)^{\beta+n-y-1}$ 是\textbf{核}，而完整的 Beta 分布还需要归一化常数：
\[
p(\theta|y) = \frac{\Gamma(\alpha+\beta+n)}{\Gamma(\alpha+y)\Gamma(\beta+n-y)} \theta^{\alpha+y-1}(1-\theta)^{\beta+n-y-1}
\]

\textbf{为什么核重要}：

\begin{enumerate}
  \item \textbf{识别分布形式}：知道核是 $\theta^a(1-\theta)^b$ 形式，就能识别这是 Beta 分布
  \item \textbf{贝叶斯推断中省略 $p(y)$}：$p(y)$ 是归一化常数，不包含关于 $\theta$ 的信息，所以可以省略
  \item \textbf{简化推导}：在推导过程中只关注核，省去繁琐的归一化常数
\end{enumerate}

★ Insight ─────────────────────────────────────
这本质上是\textbf{信息局部性}思想：归一化常数的作用只是"让面积等于1"，它本身不携带关于参数的信息。真正"有意义"的概率质量都集中在核上。
─────────────────────────────────────────────────

% === QA 记录格式 ===
% 使用以下格式记录问答：
%
% \section{问题标题}
% \label{sec:qa:LabelKey}
%
% \subsection{问题}
% 用户的原始问题
%
% \subsection{回答}
% 回答内容...
%
% 如需详细推导，添加：
% \footnote{推导见附录 \cref{sec:xxx}}

\section{学习笔记}

\textit{在这里记录你的学习心得和思考。}

% === 学习笔记格式 ===
% \section{笔记标题}
% \label{sec:notes:LabelKey}
%
% 笔记内容...


\endinput
